\documentclass[%
 reprint,
 amsmath,amssymb,
 aps,
 onecolumn,
 nofootinbib,nobibnotes,
]{revtex4-2}

% \bibliographystyle{apsrev4-1}

\usepackage{graphicx}
\usepackage{tabularx}
\usepackage{dcolumn}
\usepackage{bm}
\usepackage[dvipsnames, usenames]{xcolor}
\usepackage[hidelinks]{hyperref}
\usepackage{acronym}
\usepackage{physics}
\usepackage{orcidlink}
\usepackage{xspace}

\graphicspath{{plots}{plots/posterior_samples}{plots/prior_samples}}

\newcommand{\red}[1]{{\color{Red} #1}}
\newcommand{\jh}[1]{{\color{orange}JH: #1}}
\newcommand{\comment}[1]{}
\newcommand{\Msun}{\ensuremath{M_{\odot}}}
\newcommand{\PixelPop}{\textsc{PixelPop}\xspace}
\newcommand{\heading}[1]{\textbf{\textit{#1}}}

\newcommand{\ligo}{\affiliation{LIGO Laboratory, Massachusetts Institute of Technology, Cambridge, MA 02139, USA}}
\renewcommand{\mit}{\affiliation{Kavli Institute for Astrophysics and Space Research and Department of Physics, Massachusetts Institute of Technology, Cambridge, MA 02139, USA}}

\acrodef{GWTC-4}[GWTC-4]{fourth gravitational wave transient catalog}
\acrodef{LVK}[LVK]{LIGO---Virgo---KAGRA collaborations}



\begin{document}


\title{Sampling the Isotropic Spin Prior Conditional on $\chi_\mathrm{eff}$}

\author{Jack Heinzel\,\orcidlink{0000-0002-5794-821X}}
\email{heinzelj@mit.edu}
\ligo\mit

\date{\today}

% \begin{abstract}

% \end{abstract}

\maketitle


\section{Algorithm Description}

We wish to draw samples from
\begin{equation}
    p(a_1, a_2, c_1, c_2 \mid \chi_\mathrm{eff}, q, \mathrm{iso}),
\end{equation}
where $c_i \equiv \cos\theta_i$, the ``isotropic'' prior is
uniform in spin magnitude between 0 and $a_{\max}$ and isotropic in tilts, e.g., uniform in $\cos(\theta) \sim U(-1,1)$. The inspiral effective aligned spin parameter is defined as 
\begin{equation}
    \chi_\mathrm{eff} = \frac{a_1 c_1 + q\,a_2 c_2}{1+q}.
\end{equation}

The sampling proceeds in three stages.

\subsection{Stage 1: Marginal distribution of $\chi_i = a_i c_i$}

With $a\sim U[0,a_{\max}]$ and $c\sim U[-1,1]$ independent, let $\chi = ac$.
For fixed $a$, $p(\chi\mid a)=1/(2a)$ for $|\chi|<a$.
Marginalising over $a$:
\begin{equation}
    p(\chi \mid \mathrm{iso}) = \int_{|\chi|}^{a_{\max}}\frac{da}{2a_{\max}\, a}
    = \frac{1}{2a_{\max}}\ln\!\left(\frac{a_{\max}}{|\chi|}\right),
    \qquad |\chi|<a_{\max}.
\end{equation}
When $a_{\max}=1$ this reduces to $-\ln|\chi|/2$.

\subsection{Stage 2: Sampling $(\chi_1,\chi_2)$ given $\chi_\mathrm{eff}$ and $q$}


Bayes' theorem implies
\begin{equation}
    p(\chi_1 \mid \chi_{\rm eff}, q, \mathrm{iso}) = \frac{p(\chi_1\mid\mathrm{iso}) p(\chi_{\rm eff} \mid\chi_1, q, \mathrm{iso})}{p(\chi_{\rm eff}\mid q, \mathrm{iso})}
\end{equation}
But we can reformulate the conditional density on $p(\chi_{\rm eff} \mid\chi_1, q, \mathrm{iso})$ by realizing the only degree of freedom in $\chi_{\rm eff}$ conditional on $\chi_1$ is $\chi_2$. Changing variables to
\begin{equation}
    \chi_2 = \frac{(1+q)\chi_\mathrm{eff} - \chi_1}{q} \implies p(\chi_{\rm eff} \mid\chi_1, q, \mathrm{iso}) = \frac{1+q}{q}p(\chi_2 \mid \chi_1, q, \mathrm{iso})
\end{equation}
which introduces the Jacobian $(1+q)/q$. We also can drop the conditional dependence of $\chi_2$ on $\chi_1$ since they are independent by construction, giving the marginal
\begin{equation}
    p(\chi_1 \mid \chi_\mathrm{eff}, q, \mathrm{iso})
    = \frac{1+q}{q}
      \frac{p(\chi_1\mid\mathrm{iso})\;p(\chi_2\mid\mathrm{iso})}
           {p(\chi_\mathrm{eff}\mid q,\mathrm{iso})},
\end{equation}
where $\chi_2$ on the right-hand side is the implied value above,
and $p(\chi_\mathrm{eff}\mid q,\mathrm{iso})$ is the normalizing constant
(the convolution of the two single-spin marginals, derived in Ref.~\cite{Callister:2021gxf}). 
We sample $\chi_1$ from this one-dimensional distribution via inverse-CDF
sampling; $\chi_2$ is then determined.

\subsection{Stage 3: Sampling $(a_i, c_i)$ given $\chi_i$}

For each $\chi_i$ obtained in Stage~2, we sample $c_i$ from its conditional
distribution.  Given $\chi$ with $|\chi|<a_{\max}$, the spin magnitude must
satisfy $a = \chi/c \leq a_{\max}$, so $|c| \in [|\chi|/a_{\max},\,1]$ with the
same sign as $\chi$.  The conditional density is
\begin{equation}
    p(c\mid\chi) = \frac{1}{|c|\,\ln(a_{\max}/|\chi|)},
    \qquad |c|\in\bigl[|\chi|/a_{\max},\,1\bigr],
\end{equation}
and the corresponding CDF is
\begin{equation}
    P(c < \tau \mid \chi) = \frac{\ln\!\left(a_{\max}\,|\tau|/|\chi|\right)}
                                  {\ln\!\left(a_{\max}/|\chi|\right)},
    \qquad |\tau|\in\bigl[|\chi|/a_{\max},\,1\bigr].
\end{equation}
Inverting with a uniform draw $u\sim U[0,1]$ gives
\begin{equation}
    |c| = |\chi|^{1-u}\,a_{\max}^{\,u-1}, \qquad a = |\chi|^{u}\,a_{\max}^{\,1-u},
\end{equation}
with $c$ carrying the same sign as $\chi$.  When $a_{\max}=1$ these reduce to
$|c|=|\chi|^{1-u}$ and $a=|\chi|^u$.

\subsection{Summary}

\begin{enumerate}
    \item Draw $u_1\sim U[0,1]$ and sample $\chi_1$ by inverting the CDF of
          $p(\chi_1\mid\chi_\mathrm{eff},q,\mathrm{iso})$.
    \item Compute $\chi_2 = \bigl[(1+q)\chi_\mathrm{eff}-\chi_1\bigr]/q$.
    \item For each $\chi_i$: draw $u_i\sim U[0,1]$, set
          $c_i = \mathrm{sgn}(\chi_i)\,|\chi_i|^{1-u_i} a_{\max}^{\,u_i-1}$,
          then $a_i = |\chi_i|^{u_i} a_{\max}^{\,1-u_i}$.
\end{enumerate}

\section{Some examples}

\bibliography{draft/references} 
\end{document}
